% ===========================================================================
% Abstract
% ===========================================================================
\begin{abstract}
Google DeepMind's ``AI Agent Traps'' taxonomy identifies six categories
of environmental attacks against autonomous AI agents, yet the original
work remains purely theoretical---no open testbed exists to empirically
measure how real agents behave when deployed into adversarial
environments. We bridge this gap with \textsc{agent-traps-lab}, an
open-source framework that operationalizes all six trap categories as
22 reproducible adversarial scenarios and executes them against four
production-grade large language models (GPT-4o, Claude Sonnet~4,
Gemini~2.5~Pro, GPT-4o-mini) over Google's Agent-to-Agent (A2A)
protocol. Our experiment matrix comprises 9{,}120 individual runs
spanning three conditions---\emph{baseline} (no defenses),
\emph{hardened} (seven mitigation modules active), and \emph{ablated}
(single-mitigation removal)---each repeated 10 times for statistical
power. We report trap success rates, detection rates, escape rates,
decision drift, and cascade metrics with 95\% confidence intervals,
Wilcoxon signed-rank tests, Cohen's $d$ effect sizes, and Bonferroni
correction for multiple comparisons. Key findings include: (1)~content
injection traps via hidden CSS and HTML comments achieve the highest
baseline success rates; (2)~semantic manipulation through authority
framing is the most difficult trap category for agents to detect;
(3)~systemic multi-agent cascades produce super-additive failure modes;
(4)~the hardened mitigation suite reduces overall trap success by a
statistically significant margin with large effect sizes; and
(5)~ablation analysis reveals that input sanitization and cascade
breaking contribute the greatest marginal defense value. We release the
complete testbed, all experimental data, and paper-ready statistical
assets to enable reproducible adversarial evaluation of multi-agent
systems.
\end{abstract}
